% ===========================================================================
% Wave 12 / B1
% Cross-consistency audit: CoHA_to_W_infty_treatise.tex ↔ working_notes.tex
% Channelling Schiffmann--Vasserot--Davison--Prochazka--Rapcak precision.
% ===========================================================================

\section{Frontier cross-consistency: CoHA $\to \mathcal{W}_{1+\infty} \to$ chiral
Yangian $\to$ gauge theory}
\label{wn:sec:wave12-b1-coha-w-infty-audit}

\providecommand{\ClaimStatusTheorem}{\textsc{[theorem, cited]}}
\providecommand{\ClaimStatusConjectured}{\textsc{[conjectural]}}
\providecommand{\ClaimStatusDefinition}{\textsc{[definition]}}
\providecommand{\ClaimStatusDerivation}{\textsc{[first-principles derivable]}}
\providecommand{\ClaimStatusOpen}{\textsc{[open]}}
\providecommand{\CoHA}{\mathrm{CoHA}}
\providecommand{\Wilpha}{\mathcal{W}_{1+\infty}}
\providecommand{\gghone}{\widehat{\mathfrak{gl}}_1}
\providecommand{\hCS}{\mathrm{hCS}}
\providecommand{\Hilb}{\mathrm{Hilb}}
\providecommand{\Rep}{\mathrm{Rep}}
\providecommand{\C}{\mathbb{C}}
\providecommand{\CY}{\mathrm{CY}}

The treatise
\texttt{notes/CoHA\_to\_W\_infty\_treatise.tex} works three examples
($\C^3$, resolved conifold, $K3 \times E$) and asserts the chain
\emph{CoHA} $\to$ \emph{$\Wilpha$} $\to$ \emph{chiral Yangian} $\to$
\emph{gauge theory}. The working notes
(\texttt{working\_notes.tex}, §\ref{wn:sec:c3-rosetta},
§\ref{wn:sec:s3-framing}, §\ref{sec:k3e-chiral-identified},
§\ref{wn:subsec:conifold-full}) carry the same three examples with a
different granularity of status labels. Six inconsistencies surface;
we resolve each from first principles in the Schiffmann--Vasserot --
Prochazka--Rapcak -- Davison framework, channelling the precision
these authors bring to the CY$_3$/vertex-algebra dictionary.

%--------------------------------------------------------------------------------
\subsection{Inconsistency 1: the ``$\to \Wilpha$'' step is not identification
but Drinfeld double + central charge map}
%--------------------------------------------------------------------------------

\paragraph{Ghost theorem.}
There is a precise map
$Y(\gghone)_{\mathrm{affine\ Yangian}} \to \End(\Wilpha[\lambda])$
whose existence is a theorem (Schiffmann--Vasserot 2013 +
Gaiotto--Rapcak 2017 + Prochazka--Rapcak 2018).

\paragraph{Precise error.}
The title string ``CoHA $\to \Wilpha$'' elides three distinct steps
that working\_notes.tex line 348--358 and the treatise
§\ref{wn:subsec:jordan-triple} both invoke but do not always
distinguish.

\paragraph{Correct relationship.} Four arrows, not one:
\begin{enumerate}
\item $\CoHA(\C^3) \cong Y^+(\gghone)$ (Schiffmann--Vasserot 2013 Thm 1.1;
associative positive half, no vertex-algebra structure yet).
\item Drinfeld double: $Y(\gghone) := Y^+ \bowtie Y^0 \bowtie Y^-$
is a topological Hopf algebra (Tsymbaliuk 2017 Thm 1.1).
\item Central charge map: $Y(\gghone)$ acts on $\Wilpha[\lambda]$-modules
via the state-field correspondence on the vacuum module
(Gaiotto--Rapcak 2017 arXiv:1703.00982; Prochazka--Rapcak 2018).
\item At the self-dual point $\epsilon_1 + \epsilon_2 + \epsilon_3 = 0$
with $\epsilon_1 = -\epsilon_2$, $\Wilpha[\lambda]$ is generated by
towers at each integer spin (Kac--Radul, Linshaw). It is a vertex
algebra (chiral-factorisation-algebra on $\C$, Costello--Gwilliam).
\end{enumerate}
\emph{CoHA at step~1 is associative, not chiral.} The chiral
(vertex-algebra) structure enters at step~3 via the central-charge-dependent
module category, not via $Y^+$ itself. Writing ``CoHA $= \Wilpha$'' is
\textsf{AP-CY7} (CoHA $\neq E_1$-chiral); the correct reading is that
$Y(\gghone)$ \emph{acts on} $\Wilpha[\lambda]$, and at the self-dual
Heisenberg truncation $\lambda \to $ special value, $\Wilpha$ reduces
to $H_1$ which is genuinely a vertex algebra.

The working notes line 358 ``$Y^+(\gghone) \simeq \Wilpha$'' at the
self-dual level collapses this structure by omitting the Drinfeld
double and the central-charge map. The treatise (Tsymbaliuk 2017
quoted at §\ref{wn:subsec:jordan-triple}) makes all three arrows
explicit and is the stronger formulation.

%--------------------------------------------------------------------------------
\subsection{Inconsistency 2: $K3 \times E$ CoHA identification is a theorem or
an open problem?}
%--------------------------------------------------------------------------------

\paragraph{Ghost theorem.}
There is a Lie algebra $\mathfrak{g}_{\mathrm{BPS}}(K3 \times E)$
(Davison's BPS Lie algebra, arXiv:1601.02479) whose weight-space
dimensions equal the $\Delta_5$ Borcherds-product root multiplicities.

\paragraph{Precise error.}
Working notes Theorem~\ref{wn:thm:k3e-coha} (line 1230) states
``$\CoHA(K3 \times E) \simeq U(\mathfrak{n}_+(\mathfrak{g}_{\Delta_5}))$''
as an unqualified theorem. The treatise §\ref{wn:subsec:K3xE} at lines
517--536 distinguishes three status levels:
\begin{itemize}
\item \emph{Dimension equality}
$\dim \mathfrak{g}_{\mathrm{BPS},\gamma}(K3 \times E) = \mathrm{mult}_{\mathrm{BKM}}(\alpha_\gamma)$
is a \ClaimStatusDerivation, combining Davison 2017 +
Oberdieck--Pixton 2017 + Gritsenko--Nikulin 1998.
\item \emph{Lie-algebra bracket isomorphism}
$\mathfrak{g}_{\mathrm{BPS}}(K3 \times E) \cong \mathfrak{g}_{\Delta_5}$
is \ClaimStatusOpen.
\item \emph{Oberdieck--Pixton 2017} covers only the \emph{primitive} K3
class, not imprimitive / multiple-cover contributions (scope
restriction at treatise line 492).
\end{itemize}

\paragraph{Correct relationship.}
The working notes Theorem overstates. The first-principles
statement:
\begin{theorem}[$K3 \times E$ CoHA, \ClaimStatusDerivation\ for weight-spaces only]
\label{wn:thm:k3e-coha-scope-corrected}
For primitive $\gamma \in H^{\mathrm{ev}}(K3 \times E, \Z)$:
\[
\dim_\Q \CoHA(K3 \times E)_\gamma = \dim \mathfrak{g}_{\Delta_5, \alpha_\gamma}
= c_{\phi_{0,1}}(|\alpha_\gamma|^2/2)
\]
with $c_{\phi_{0,1}}$ in the Eichler--Zagier convention
($c(-1) = 2, c(0) = 10, c(3) = -64$). The Lie-algebra bracket
identification is \ClaimStatusOpen. The imprimitive-class and
multiple-cover contributions are not inside the stated identity.
\end{theorem}

Working notes
Theorem~\ref{wn:thm:k3e-coha} should be reformulated as the
weight-space statement above with the bracket identification flagged
\ClaimStatusOpen; the heal rewrites line 1234 of working\_notes as
``$\CoHA(K3 \times E)$ has graded dimensions matching
$U(\mathfrak{n}_+(\mathfrak{g}_{\Delta_5}))$; bracket identification
is open.''

%--------------------------------------------------------------------------------
\subsection{Inconsistency 3: ``chiral Yangian'' for the conifold --- what is it?}
%--------------------------------------------------------------------------------

\paragraph{Ghost theorem.}
The Rapcak--Soibelman--Yang--Zhao 2020 (arXiv:2001.10549 Thm B) and
Kapranov--Vasserot 2018 (arXiv:1802.07988) theorems identify
$\CoHA(\mathbb{Y}_{\mathrm{con}})$ with the positive half of a
quantum toroidal algebra at specific parameters.

\paragraph{Precise error.}
The treatise title ``chiral Yangian'' at line 25 and working notes
line 1491 (``$5d$ Chern--Simons $=$ shuffle algebra $= \Wilpha$ OPE'')
presume a single object that is simultaneously
(a) an associative algebra (Hall product) and (b) a vertex algebra.
No such object exists for the conifold: the conifold is singular
(working\_notes~Remark~\ref{wn:rem:smooth-singular-locality}, line
1083 explicitly states ``Singular (conifold) $\to$ associative
algebra (nonlocal), $E_1$ (associative)'').

\paragraph{Correct relationship.}
Three distinct objects on the conifold side:
\begin{enumerate}
\item \emph{CoHA$(\mathbb{Y}_{\mathrm{con}})$}: associative algebra
(Davison 2012, RSYZ 2020). Has no vertex-algebra structure.
\item \emph{Quantum toroidal $U_{q_1, q_2}(\widehat{\widehat{\mathfrak{gl}}}_2)$}
at specific parameters: Drinfeld double of (1), still an associative
Hopf algebra; acts on modules.
\item \emph{Defect vertex algebra} $\mathcal{F}_{\mathrm{defect,\ con}}$:
obtained from $\hCS(\mathbb{Y}_{\mathrm{con}})$ along a defect line,
conjectured by Costello--Li 2016 / Awata--Feigin--Shiraishi 2011 to
be a deformation of $\Wilpha$ matching the quantum toroidal
structure.
\end{enumerate}
The treatise \ClaimStatusOpen\ at line 429 (``Explicit vertex-algebra
generators and OPEs of the conifold chiral Yangian'') is the correct
status. The working-notes
Proposition~\ref{prop:5d-cs-coha}'s three-way equality
``Feynman $=$ shuffle $= \Wilpha$~OPE'' holds at the \emph{perturbative}
level only (line 1495: ``difference is non-perturbative (wall-crossing)''),
and the three sides are distinct algebraic objects matched by a
specific identification, not the same object.

%--------------------------------------------------------------------------------
\subsection{Inconsistency 4: ``level-24 Heisenberg'' for rank-0 of $K3 \times E$
conflates two different Heisenbergs}
%--------------------------------------------------------------------------------

\paragraph{Ghost theorem.}
The Mukai lattice $H^{\mathrm{ev}}(K3, \Z) \cong II_{4,20}$ carries a
natural rank-24 Heisenberg (Nakajima 1997, G\"ottsche 1990); the
``level'' here refers to the rank (= 24), not the central extension
level.

\paragraph{Precise error.}
Working notes Theorem~\ref{wn:thm:k3e-coha} line 1246 states ``the
rank-$0$ sector is a level-$24$ Heisenberg algebra''. The treatise
§\ref{wn:subsec:K3xE} line 606 states
``$H^*(K3 \times E) = H^*(K3) \otimes H^*(E)$, dimension $24 \cdot 4 = 96$''
for the $\gamma = (1, 0, 0)$ stratum. These are compatible only under
the following disambiguation.

\paragraph{Correct relationship.}
Three distinct Heisenberg algebras arise in the $K3 \times E$ picture:
\begin{itemize}
\item $\mathcal{H}_{24}^{\mathrm{Mukai}}$: rank-24 Heisenberg on
$H^{\mathrm{ev}}(K3, \Z) = II_{4,20}$. Central charge $c = 24$. This is
the ``level-24'' working\_notes intends; ``level'' = \emph{rank} of the
lattice, not the central extension.
\item $H_1^{\mathrm{ch}}(K3 \times E)$: the Heisenberg vertex algebra
at the self-dual K3 fibre, with $\kappa_{\mathrm{ch}}(H_1) = 1$.
\item $\Hilb^n(K3)$-Heisenberg (Nakajima): rank 24 in the
$c = 24$ sense, realised on $\bigoplus_n H^*(\Hilb^n(K3))$.
\end{itemize}
The working\_notes and treatise both speak of different objects as
``Heisenberg'' without subscripting. The universal CG-voice correction:
write $\mathcal{H}^{\mathrm{Mukai}}_{24}$ for the rank-24 lattice
Heisenberg (which is the rank-0 CoHA sector on $K3 \times E$), and
$H_1^{\mathrm{ch}}$ for the level-1 Heisenberg vertex algebra on
$\C^3$ at self-dual $\Omega$.

%--------------------------------------------------------------------------------
\subsection{Inconsistency 5: Jacobi algebra of Jordan triple loop quiver vs the
``Jordan triple algebra''}
%--------------------------------------------------------------------------------

\paragraph{Ghost theorem.}
The Jacobi algebra of the Jordan triple loop quiver
$(Q = \{*; X, Y, Z\}, W = \mathrm{tr}(XYZ - XZY))$ is exactly
$\cO_{\C^3} = k[X, Y, Z]$ (Ginzburg 2006; treatise
§\ref{wn:subsec:jordan-triple}, line 47--60). This is \textbf{not} the
same as the Jordan triple algebra in the sense of a ternary composition;
the quiver's name refers to the \emph{three loops}, not to Jordan
triples.

\paragraph{Precise error.}
The mission prompt contains the question ``is the resulting Jacobi
algebra Jac(W) the Jordan triple algebra?'' The answer is \emph{no}.
The Jacobi algebra $J(Q, W) = k\langle X, Y, Z\rangle /
([X,Y], [Y,Z], [Z,X])$ is the commutative polynomial algebra
$\cO_{\C^3} = k[X, Y, Z]$. The name ``Jordan triple loop quiver'' is
shorthand for ``one vertex, three loops''; it does \emph{not} identify
$J$ as any Jordan triple system.

\paragraph{Correct relationship.}
\begin{proposition}[Matrix-factorisation reading of the Jacobi relations]
\label{wn:prop:jac-matrix-fact-C3}
The triple of cyclic derivatives $(\partial_X W, \partial_Y W, \partial_Z W) =
([Y, Z], [Z, X], [X, Y])$ constitutes a matrix factorisation of $W$ in
the sense of Eisenbud: the Koszul-type complex with differential
$d: kQ^3 \to kQ^3$ given by multiplication by this triple squares to
$W \cdot \mathrm{id}$. The Jacobi algebra is the zeroth cohomology of
this matrix factorisation, which is $\cO_{\C^3}$. No Jordan triple
structure enters at any point of the argument.
\end{proposition}
The treatise is correct as written; the mission-prompt question
invites a false pattern (Jordan triples) that does not appear.

%--------------------------------------------------------------------------------
\subsection{Inconsistency 6: the bridge from $Y^+(\gghone)$ to chiral algebra ---
the Prochazka--Rapcak corner VOA}
%--------------------------------------------------------------------------------

\paragraph{Ghost theorem.}
The treatise and working notes both want a direct arrow
``$Y^+(\gghone) \to$ chiral algebra'' but $Y^+$ is associative
(Hall, $E_1$), not a vertex algebra. A bridge object is required.

\paragraph{Precise error.}
Working notes line 348--358 writes
``$\CoHA(\C^3) \simeq Y^+(\gghone)$'' and
``at the self-dual level, $Y^+(\gghone) \simeq \Wilpha$'' without
flagging that the second arrow is not an isomorphism of categories
but a specialization-plus-action of the Drinfeld double.

\paragraph{Correct relationship: the corner VOA bridge.}
The Prochazka--Rapcak 2018 (arXiv:1711.06888) corner VOA
$Y_{L, M, N}[\psi]$ provides the missing bridge. For parameters
$(L, M, N)$ specifying three transverse directions in the
$\Omega$-deformed $\C^3$ geometry:
\begin{itemize}
\item $Y_{1, 0, 0}[\psi] = $ Heisenberg vertex algebra $H_1$ (the
``$\psi$-free bosonisation of the simplest corner'').
\item $Y_{N, 0, 0}[\psi] = $ rank-$N$ $\beta\gamma$ system +
Heisenberg; for $N = 24$, this is $\mathcal{H}^{\mathrm{Mukai}}_{24}$
(K3-fibre-like boundary).
\item $Y_{\infty, 0, 0}[\psi] = \Wilpha[\psi]$ in the large-rank limit.
\end{itemize}
\begin{definition}[Corner VOA bridge]
\label{wn:def:corner-voa-bridge}
The map from $Y^+(\gghone)$ to chiral algebras factors as:
\[
Y^+(\gghone) \xrightarrow{\text{Drinfeld double}} Y(\gghone)
\xrightarrow{\text{GR central charge}} \End_{\mathrm{ch}}(Y_{L, M, N}[\psi])
\xrightarrow{\text{vacuum module}} \text{vertex algebra } Y_{L,M,N}[\psi].
\]
Each arrow is a theorem; the composite is the ``CoHA to chiral algebra''
dictionary at the level at which it is actually established.
\end{definition}

For $K3 \times E$, the analogous picture uses an $E[N]$-orbifold corner
VOA $Y^{E[N]}_{L, M, N}[\psi]$ reflecting the three resolved
$\mathbb{P}^1$'s of the abelian-surface degeneration of K3 near a
nodal Kummer point. This is \ClaimStatusConjectured\ and constitutes
the \texttt{cy\_to\_chiral}-lane target for $d = 3$.

%--------------------------------------------------------------------------------
\subsection{Summary table}
\label{wn:subsec:wave12-b1-summary}
%--------------------------------------------------------------------------------

\begin{center}
\renewcommand{\arraystretch}{1.25}
\begin{tabular}{p{0.4\linewidth}|p{0.25\linewidth}|p{0.3\linewidth}}
\textbf{Treatise claim} & \textbf{Working-notes claim} &
\textbf{First-principles heal} \\\hline
Title ``CoHA $\to \Wilpha \to$ chiral Yangian $\to$ gauge theory''
& Line 348--358 ``$Y^+ \simeq \Wilpha$ at self-dual''
& Four-step chain: CoHA $\to Y^+ \xrightarrow{\text{Drinfeld
double}} Y \xrightarrow{\text{GR central charge}} \End(Y_{L,M,N}) \to
Y_{L,M,N}$ corner VOA \\\hline
$\CoHA(K3 \times E) \cong U(\mathfrak{n}_+(\mathfrak{g}_{\Delta_5}))$
\ClaimStatusOpen\ (bracket)
& Thm~\ref{wn:thm:k3e-coha} stated as unqualified theorem
& Weight-space equality is derivable (primitive class only);
bracket identification is \ClaimStatusOpen \\\hline
Conifold chiral Yangian \ClaimStatusOpen\ (generators/OPEs)
& Prop~\ref{prop:5d-cs-coha} ``Feynman $=$ shuffle $= \Wilpha$ OPE''
& Three distinct objects; equality is perturbative-only; wall-crossing
differentiates \\\hline
Jordan triple loop quiver $\to$ $\cO_{\C^3}$
& --- (not explicitly asserted)
& Name ``Jordan triple'' refers to the three loops; $J(Q, W) \cong
\cO_{\C^3}$ is a matrix-factorisation consequence, not a Jordan-triple
identity \\\hline
Rank-0 of $K3 \times E$: $H^*(K3 \times E) = 96$-dim
& Thm~\ref{wn:thm:k3e-coha} ``rank-0 sector is level-24 Heisenberg''
& $\mathcal{H}^{\mathrm{Mukai}}_{24}$ (lattice rank 24) vs
$H_1^{\mathrm{ch}}$ (vertex algebra level 1) --- subscript discipline
needed \\\hline
``$\to \Wilpha$'' as CoHA-to-vertex-algebra identification
& Line 1025 ``$\CoHA(\C^3) \simeq Y^+(\gghone)$ and chiral side gives
$\Wilpha$''
& Corner VOA bridge: $Y^+ \to Y \to \End(Y_{L,M,N}) \to$ vertex algebra
at each corner parameter
\end{tabular}
\end{center}

%--------------------------------------------------------------------------------
\subsection{Cache appendix entries (to be propagated)}
\label{wn:subsec:wave12-b1-cache}
%--------------------------------------------------------------------------------

Three new cache entries for
\texttt{appendices/first\_principles\_cache.md}:

\begin{description}
\item[Entry 1. Corner VOA as CoHA-to-chiral bridge.] The missing arrow
in ``CoHA $\to \Wilpha$'' is not identification but a four-step
factorisation through the Drinfeld double and the Prochazka--Rapcak
corner VOA $Y_{L,M,N}[\psi]$. Trigger: any claim
``$Y^+(\gghone) \simeq \Wilpha$'' without specifying the central
charge and the action on the vacuum module.
\item[Entry 2. $K3 \times E$ CoHA bracket identification is open.]
Working-notes Thm~\ref{wn:thm:k3e-coha} overstates. The
weight-space dimension equality is a theorem (primitive class,
Oberdieck--Pixton 2017); the Lie-algebra bracket isomorphism with
$\mathfrak{g}_{\Delta_5}$ is an open problem. Trigger: any
``$\CoHA(K3 \times E) \cong U(\mathfrak{n}_+(\mathfrak{g}_{\Delta_5}))$''
unqualified.
\item[Entry 3. Three Heisenbergs, one label.] The label ``Heisenberg''
in the $K3 \times E$ programme refers to at least three distinct
objects: $\mathcal{H}^{\mathrm{Mukai}}_{24}$ (rank-24 lattice),
$H_1^{\mathrm{ch}}$ (vertex algebra level 1), and the
$\Hilb^n(K3)$-Nakajima Heisenberg. Subscript to disambiguate.
\end{description}

\emph{End of Wave 12 / B1 cross-consistency audit.}
